\documentclass[11pt]{article}
\usepackage[margin=0.9in]{geometry}
\usepackage{amsmath,amssymb,amsthm,booktabs,longtable,array}
\usepackage[T1]{fontenc}
\usepackage{lmodern,microtype}
\usepackage[colorlinks=true,linkcolor=blue,urlcolor=blue]{hyperref}
\newtheorem{theorem}{Theorem}[section]
\newtheorem{lemma}[theorem]{Lemma}
\newtheorem{proposition}[theorem]{Proposition}
\newtheorem{corollary}[theorem]{Corollary}
\theoremstyle{remark}
\newtheorem{remark}[theorem]{Remark}
\newcommand{\R}{\mathbb R}
\newcommand{\C}{\mathbb C}
\newcommand{\E}{\mathbb E}
\newcommand{\norm}[1]{\left\|#1\right\|}
\newcommand{\abs}[1]{\left|#1\right|}
\newcommand{\Herm}{\operatorname{Herm}}
\newcommand{\tr}{\operatorname{tr}}
\newcommand{\Ad}{\operatorname{Ad}}
\title{Toward a four-dimensional Yang--Mills construction\\
\large Notebook 05: the flat gauge bridge and cutoff interactions}
\author{Research working document}
\date{16 October 2026}
\setlength{\emergencystretch}{5em}
\begin{document}
\maketitle

\begin{abstract}
This note proves two bridges needed to apply the earlier Gaussian
estimates to the lattice construction. First, it identifies the actual
averaged-path gauge penalty, distinguishes the rooted and block-mean
constraints, and transfers a flat axial gap quantitatively to the
exponential gauge. Second, physical-field analyticity gives the correct
enlarged-cutoff bounds for the interaction, including the linear
Jacobian term. A complete Gaussian cluster argument is given for
interactions within disjoint fixed-dimensional cells, with an explicit
bound on the change of the logarithm. The actual Yang--Mills interaction
joins overlapping cells; its integration theorem is still a separate,
unverified interface. No continuum theory or mass gap is asserted.
\end{abstract}

\noindent\textbf{Status and sources.}
These are internally checked finite-dimensional and conditional
estimates, without external mathematical review or a novelty claim.
Raw I and Raw II refer to
\href{https://github.com/ethanwsussman/Yang--Mills-LLM-experiment/tree/e392ef45013a581664bab76b90867170b30fe10e}
{the pinned deposited account}.
The model formulas are compared with the original articles
\cite{balaban99,balaban109,balaban116}.
Notebook 03 supplies the Gaussian forest and logarithm estimates;
Notebook 04 supplies the stronger local-precision support algebra.
Their hypotheses are retained whenever invoked below.

\section{The actual flat gauge penalty and its constraint}

Work on a unit lattice partitioned into blocks of side $L$, with
$n=L^d$ sites and a distinguished center $y$ in each block $B(y)$.
Count one positive orientation per bond. Lie-algebra coordinates are
Hermitian, with invariant norm $|X|^2=\tr(X^2)$ in a fixed normalization.
The following statements use the ordinary block tiling with distinct
interfaces. A compatible Dirichlet convention takes a region made of
complete blocks, retains bonds with at least one endpoint in it, sets
fields to zero on bonds with both endpoints outside, and sets gauge
parameters to zero outside. Blockwise mean subtraction then preserves
the allowed site space and the boundary pivots remain available.
This is the incident-bond convention in Raw II, line 9033 and the
positivity hypothesis beginning at line 9743. Freezing all crossing
bonds would require a different argument. Periodic seams must remain
distinct geometric edges, even with only two blocks in a direction.

For each $x\in B(y)$, let $\mathcal P(y,x)$ be the prescribed family
of monotone shortest paths. Define the averaged path integral
\[
 (TB)(y,x)=\mathop{\rm average}_{\Gamma\in\mathcal P(y,x)}
                 \sum_{b\in\Gamma}\epsilon_b B_b,
 \qquad (TB)(y,y)=0.
\]
Let $T_0$ denote a single rooted-tree path integral, choosing a tree
whose root-to-site paths have length at most $h=d(L-1)$. Write
$D\lambda(v,w)=\lambda(w)-\lambda(v)$ and
$d\lambda(y,x)=\lambda(x)-\lambda(y)$. Telescoping gives
\begin{equation}\label{eq:telescoping}
 TD\lambda=T_0D\lambda=d\lambda.
\end{equation}
On rooted site fields $\lambda(y)=0$, this is $TD\lambda=\lambda$.
The useful elementary bounds are
\begin{equation}\label{eq:TD-bounds}
 \norm D_{2\to2}\le d_0:=2\sqrt d,\qquad
 \norm T_{2\to2},\norm{T_0}_{2\to2}
       \le t_0:=\sqrt{h(n-1)}.
\end{equation}
For $D$, sum $|\lambda(v)-\lambda(w)|^2\le
2|\lambda(v)|^2+2|\lambda(w)|^2$. For $T$, each absolute
coefficient row sums to at most $h$, and each column to at most
$n-1$; the Schur test proves the assertion.

The gauge fixing in CMP109, equations (0.11), (0.14)--(0.16), uses
these averaged contours. If its group variable is
$\Omega(t)=\exp(i[tTB+O(t^2)])$, then
\[
 1-\Re\tr\Omega(t)=\tfrac12t^2|TB|^2+O(t^3).
\]
Thus its actual quadratic penalty, with the stated normalization,
is $\tfrac12\norm{TB}^2$. The second-order term in the logarithm
does not contribute to this coefficient. This matches CMP109 (2.5).
For a tree-axial field $B_0$, generally $TB_0\ne0$. Hence on
$B_0+D\lambda$ the penalty is
$\tfrac12\norm{TB_0+\lambda}^2$, requiring a shear of the
gauge coordinate before it becomes diagonal.

\subsection{Rooted and block-mean constraints are different}

For a coarse bond $c=(y,y+Le_\mu)$, set
\[
 (Q_{\rm st}B)(c)=L^{-d-1}\sum_{x\in B(y)}
                         \int_{[x,x+Le_\mu]} B,
\]
\[
 (aTB)(y)=L^{-d}\sum_{x\in B(y)}(TB)(y,x),\qquad
 (D_1f)(c)=L^{-1}[f(y+Le_\mu)-f(y)].
\]
Linearizing the rooted contour average of CMP109 (0.12) gives
\begin{equation}\label{eq:Q-dictionary}
 Q^\flat=Q_{\rm st}-D_1aT,\qquad
 Q_1:=Q^\flat+D_1aT=Q_{\rm st}.
\end{equation}
Indeed the contour integral is the initial rooted path, plus the
translated straight string, minus the final rooted path and the
reference coarse string. Multiplication by the reference transporter
cancels the last term. Averaging and dividing by $L$ gives
\eqref{eq:Q-dictionary}. In particular,
\begin{equation}\label{eq:Q-gauge}
 Q^\flat D\lambda=D_1(\lambda\hbox{ at centers}),\qquad
 Q_1D\lambda=D_1(\hbox{block mean of }\lambda).
\end{equation}
Rooted gauge directions preserve $\ker Q^\flat$; block-mean-zero
directions preserve $\ker Q_1$. The two constraints agree on $\ker T$.

For $B\in\ker Q^\flat$, take the rooted field $\xi=TB$ and put
\begin{equation}\label{eq:B-T}
 B_T=B-D\xi.
\end{equation}
Then $TB_T=0$, $Q^\flat B_T=Q_1B_T=0$. Conversely any
$B_T\in\ker Q_1\cap\ker T$ and any rooted $\xi$ give a unique
$B=B_T+D\xi\in\ker Q^\flat$ with $TB=\xi$.

To compare with a tree slice inside the correct constraint, use the
block-mean-zero site space $N$. The inverse of $d:N\to\{f(y,y)=0\}$ is
\[
 (d^{-1}f)(x)=f(y,x)-L^{-d}\sum_{z\in B(y)}f(y,z),
 \qquad \norm{d^{-1}f}_2\le\norm f_2.
\]
The norm inequality is orthogonal subtraction of the block mean.
The mutually inverse maps between the two axial slices are
\begin{equation}\label{eq:slice-maps}
 B_0=B_T-Dd^{-1}T_0B_T,\qquad
 B_T=B_0-Dd^{-1}TB_0.
\end{equation}
They preserve $Q_1$ by \eqref{eq:Q-gauge}; telescoping verifies
the slice conditions and both inverse identities. In particular
\begin{equation}\label{eq:slice-norm}
 \norm{B_T}\le(1+d_0t_0)\norm{B_0}.
\end{equation}

\section{A quantitative flat coercivity transfer}

\begin{theorem}\label{thm:gauge-gap}
Let $q$ be a real symmetric quadratic form on bond fields with
$q(B+D\lambda)=q(B)$ for the relevant gauge directions. Suppose
\[
 q(B_0)\ge\gamma_{\rm ax}\norm{B_0}^2
 \quad\hbox{on }\ker Q_1\cap\ker T_0.
\]
Then, on $\ker Q^\flat$,
\begin{equation}\label{eq:exponential-gap}
 q(B)+\norm{TB}^2\ge\gamma_{\rm exp}\norm B^2,
 \qquad
 \gamma_{\rm exp}=
 \left(\frac{(1+d_0t_0)^2}{\gamma_{\rm ax}}+d_0^2\right)^{-1}.
\end{equation}
If the penalty is $a\norm{TB}^2$ with $a>0$, replace $d_0^2$
by $d_0^2/a$ in the denominator.
\end{theorem}

\begin{proof}
Gauge invariance, \eqref{eq:slice-maps}, and
\eqref{eq:slice-norm} give
$q(B_T)\ge\gamma_T\norm{B_T}^2$, where
$\gamma_T=\gamma_{\rm ax}/(1+d_0t_0)^2$.
The exact decomposition \eqref{eq:B-T} yields
\[
 q(B)+\norm{TB}^2=q(B_T)+\norm\xi^2
 \ge\gamma_T\norm{B_T}^2+\norm\xi^2.
\]
Weighted Cauchy--Schwarz applied to
$\norm B\le\norm{B_T}+d_0\norm\xi$ gives
\eqref{eq:exponential-gap}. The version with $a$ is identical.
\end{proof}

The axial input here concerns the full constrained bond norm. A
statement in independent coordinates $B_0=K_{\rm ax}z$ with
$q(K_{\rm ax}z)\ge\gamma_{\rm free}\norm z^2$ only gives
$\gamma_{\rm ax}=\gamma_{\rm free}/\norm{K_{\rm ax}}^2$.
This conversion is essential.

\subsection{Explicit flat pivot-coordinate bounds}

Choose one pivot $b_0(c)$ for each coarse bond: the unit bond crossing
its block interface on the line through the two centers. No rooted
contour in $T$ uses this interblock bond. Exactly $L$ starting sites
in the straight-string sum use it, so its coefficient in $Q^\flat$
is $L^{-d}$. No other coarse equation uses the pivot. Therefore,
after splitting the remaining coordinates $z$,
\[
 B_{\rm pivot}=-L^dQ_fz,\qquad
 \norm{K_\flat}^2\le1+L^{2d}\norm{Q^\flat}^2.
\]
Put $r=1+2h/L$. The absolute row sum of $Q^\flat$ is at most $r$.
An internal bond can enter contour terms only through its own block,
which has $2d$ incident coarse bonds; straight strings add at most
two equations. Thus the absolute column sum is at most $(2d+2)r$,
and
\begin{equation}\label{eq:pivot-bounds}
 \norm{K_\flat}^2\le1+L^{2d}(2d+2)r^2,
 \qquad \norm{K_{\rm ax}}^2\le1+2L^{2d}.
\end{equation}
For the second bound, tree bonds vanish and $Q_1=Q_{\rm st}$ has
row sum one and column sum at most two; the same pivots remain free
of the block trees. These conservative constants are independent of
scale and volume for this tiling. A different pivot scheme or a
boundary that identifies interfaces needs its own coefficient matrix.

If also $q(B_0)\le M_q\norm{B_0}^2$ on the axial slice, the
projections above imply
\[
 q(B)+\norm{TB}^2
 \le[M_q(1+d_0t_0)^4+t_0^2]\norm B^2.
\]
After $B=K_\flat z$, the lower bound remains
$\gamma_{\rm exp}\norm z^2$ since $z$ is a coordinate restriction
of $B$, and the upper bound gains $\norm{K_\flat}^2$.
Both bounds matter for a uniformly accretive complex covariance.

\subsection{Why this applies at the flat reference, and its limit}

For the actual gauge-invariant effective action
$f(V)=\mathcal A(U_k(V))$, the trivial datum is stationary: the
Wilson action is nonnegative and vanishes at the trivial field.
The Hessian therefore annihilates infinitesimal gauge directions.
The gradient-times-second-chart-derivative correction also vanishes
there. Consequently the flat action Hessian has exactly the invariance
needed in Theorem~\ref{thm:gauge-gap}, and the actual gauge penalty is
the one derived above. This identifies a bridge from a uniform axial
estimate, not an independent proof of that axial estimate for every
effective scale.

The algebraic distinctions already occur in Raw II, lines 9045--9285;
the deposited account's later quantitative covariance statement has
additional hypotheses. The remaining model inputs include the primary
axial gap and upper bound in the matched coordinates, and uniform
weighted perturbation estimates for the actual complex independent-current
family. A change of coordinates in a coercivity proof does not authorize
a change of the real sharp-cutoff integration domain.

\section{The current correction to gauge invariance}

Use $U'(b)=\exp(i\eta B_b)U_b$ and the convention
$U^g(v,w)=g(v)^{-1}U(v,w)g(w)$. Expand the actual action as
\[
 f_U(B)=f(U)+\langle J_U,B\rangle+\tfrac12H_U(B,B)+O(\norm B^3).
\]
For $g_t=\exp(it\lambda)$, write $a=\lambda(v)$,
$b'=\Ad_{U(v,w)}\lambda(w)$ and
$D_U\lambda=(b'-a)/\eta$. The Baker--Campbell--Hausdorff expansion
of the transformed two-parameter field gives the mixed coordinate
\[
 sB+tD_U\lambda+\tfrac{i}{2}st[B,a+b']
       +O(s^2t+st^2)+\hbox{terms independent of }s.
\]
Gauge invariance therefore proves the exact Ward identity
\begin{equation}\label{eq:ward}
 H_U(B,D_U\lambda)=-\tfrac{i}{2}
                 \langle J_U,[B,a+b']\rangle.
\end{equation}
It also gives $\langle J_U,D_U\lambda\rangle=0$.
The pure-gauge second coordinate derivative is
$-(i/\eta)[a,b']$, agreeing with
$H_U(D_U\lambda,D_U\lambda)=(i/\eta)
\langle J_U,[a,b']\rangle$ and checking the sign and scale.

For unitary $U$, choose $c_{\rm br}$ with
$|[X,Y]|\le c_{\rm br}|X||Y|$, and put
$j_* =\sup_b|J_U(b)|$. Since the square sum of
$|\lambda(v)|+|\lambda(w)|$ is at most $4d\norm\lambda^2$,
\[
 |H_U(B,D_U\lambda)|
 \le c_{\rm br}\sqrt d\,j_*\norm B\norm\lambda.
\]
At unit spacing, the change of the quadratic form obeys
\begin{align*}
 |H_U(B+D_U\lambda,B+D_U\lambda)-H_U(B,B)|
 \le{}&2c_{\rm br}\sqrt d\,j_*\norm B\norm\lambda\\
 &+2c_{\rm br}d\,j_*\norm\lambda^2.
\end{align*}
Thus no derivative of the current is required for this correction.
Complex backgrounds require explicit bounds on their adjoint transports.
An arbitrary independent current is not automatically the derivative
of an invariant action; its analytic extension must be checked through
the actual operator formula. Raw II, lines 9380--9560, already gives
the corresponding physical Ward calculation with the opposite gauge
convention.

There is an even sharper physical-fiber statement. If the averaging
level set is preserved by rooted gauges and the field is a critical
point of the action restricted to that level set, the complete Hessian
in a level-set chart annihilates rooted gauge directions. The mixed
derivative of gauge invariance equals the Hessian pairing plus a first
derivative applied to the mixed chart term; the latter is zero by
constrained criticality. This includes the constraint-curvature/current
correction. It supplies no analogous identity for a freely prescribed
current and no axial gap by itself.

\section{Interaction smallness at the larger real cutoff}

\begin{lemma}[Physical-field Schwarz estimate]\label{lem:schwarz}
Let $F_Y(A,z)$ be holomorphic on the balanced physical-field ball
$\norm z_{\infty,Y}<r_0$, uniformly over its stated background domain.
Assume $|F_Y|\le a_Y$ and $D_z^jF_Y(A,0)=0$ for $j=0,1,2$.
Then
\[
 |F_Y(A,z)|\le a_Yr_0^{-3}\norm z_{\infty,Y}^3.
\]
For real $B$, $\norm B_{\infty,Y}\le R$, and $gR<r_0$, the
action remainder $W_Y=g^{-2}F_Y(A,gB)$ satisfies
\begin{equation}\label{eq:cubic-stability}
 |W_Y|\le\frac{a_Y}{r_0^3}gR\norm B_{2,Y}^2,
 \qquad \sup_{\rm cutoff}|W_Y|\le\frac{a_Y}{r_0^3}gR^3.
\end{equation}
\end{lemma}

\begin{proof}
For $z\ne0$, apply Schwarz's lemma three times to
$w\mapsto F_Y(A,r_0wz/\norm z_\infty)/a_Y$ on the unit disk.
The zero cases and non-strict limits are immediate. Substitute $z=gB$
and use $\norm B_\infty^3\le R\norm B_2^2$.
\end{proof}

If $\sup_b\sum_{Y\ni b}a_Y\le a_*$, summing gives
\[
 \sum_Y|W_Y|\le a_*gRr_0^{-3}\sum_b|B_b|^2.
\]
This can be absorbed into a fixed real Gaussian quadratic exponent
when $gR$ is sufficiently small. More generally, a term
$g^{-r}F_Y(gB)$ vanishing to order $m$ has cutoff norm at most
$a_Yg^{m-r}R^m/r_0^m$. A local norm
$\sup_b\sum_{Y\ni b}e^{a|Y|+\kappa\ell(Y)}a_Y$ transfers with
the same factor. An extensive bound on the sum of the functions does
not imply such a local norm. A physical argument such as $gCB$ or a
nonlinear minimizer response needs its own uniform bound and local
support expansion before this lemma can be applied.

\subsection{The full exponent is not cubic}

The full fluctuation potential in CMP109, pages 266--270, and Raw II,
lines 73910--73940, also includes change-of-variables and Haar factors.
Its schematic orders in the physical argument $z=gCB$ are:
\begin{center}
\begin{tabular}{@{}lll@{}}
\toprule
Term & First order & Cutoff scale\\
\midrule
Logarithm of the Jacobian & 1 & $gR$\\
Haar log density & 2 & $g^2R^2$\\
Action, gauge, and current cubic remainders & 3 with $g^{-2}$ & $gR^3$\\
Quadratic term in the quadratic constraint correction & 4 with $g^{-2}$ & $g^2R^4$\\
\bottomrule
\end{tabular}
\end{center}
These orders use the stated cancellations and do not establish their
localized coefficient bounds. Previous effective-action increments
are separate induction inputs. For $R\ge1$ and $gR\le1$, all displayed
scales are bounded by a constant times $gR^3$. Thus such coefficient
bounds imply $u_Y:=\sup_{\rm cutoff}|V_Y|\le C_YgR^3$.
For $u_Y\le1$,
\begin{equation}\label{eq:interaction-v}
 v_Y:=\sup_{\rm cutoff}|e^{V_Y}-1|\le e^{u_Y}-1\le2u_Y.
\end{equation}

A nonzero linear term cannot obey $|V_Y(B)|\le CgR\norm B^2$
near zero. An additive constant is necessary, for example
$Cg\norm B\le(Cg/2)(1+\norm B^2)$. The above orders allow a
safe envelope $\widetilde C_Yg(1+R)(1+\norm B_{2,Y}^2)$,
with summable localized constants. Odd Gaussian cancellations do not
repair a false pointwise quadratic bound.

\section{A complete interaction theorem for disjoint cells}

Let $\Lambda$ be a finite set of cells of fixed real dimension $D$,
with metric $d$. Suppose their holomorphic complex symmetric covariance
obeys
\[
 \Herm C(A)\ge\lambda I,\qquad \norm{C(A)}\le\Lambda_0,
 \qquad \norm{C_{ij}(A)}\le J_{ij}=J_{ji},
\]
\[
 A_\kappa:=\sup_i\sum_{j\ne i}e^{\kappa d(i,j)}J_{ij}<\infty.
\]
Each interaction $V_i(A,x_i)$ reads only its own cell's fluctuation.
It is measurable in real $x_i$, holomorphic in the background on the
fixed domain, and obeys
\[
 \sup_{A,|x_i|\le R}|e^{V_i(A,x_i)}-1|\le v.
\]
Put $\chi_i=\mathbf1_{\{|x_i|\le R\}}$ and define globally
\[
 q_i=-\mathbf1_{\{|x_i|>R\}}+\chi_i(e^{V_i}-1).
\]
The values of $V_i$ outside the cutoff are unnecessary. Then
\[
 Z_V(A)=\E_{C(A)}\prod_i(1+q_i)
       =\E_{C(A)}\left(\prod_i\chi_i\right)e^{\sum_iV_i}.
\]
The complex Gaussian uses the normalized analytic determinant branch
on the accretive cone. Let
\begin{gather*}
 \delta=\lambda/2,\quad V_*=8\Lambda_0^2/\lambda,
 \quad K_G=(2V_*/\lambda)^{D/2},
 \quad p_R=2^{D/2}e^{-R^2/(4V_*)},\\
 \beta=K_G(p_R+v),\quad B_0=8e^2A_\kappa/\delta,
 \quad \beta_*=[2e^2\max(1,B_0)]^{-1}.
\end{gather*}
Here $\ell(X)$ denotes minimal spanning-tree length in the cell metric.

\begin{theorem}\label{thm:cell-interaction}
If $\beta\le\beta_*$, there is a convergent analytic logarithm
\[
 \log Z_V(A)=\sum_{\varnothing\ne X\subset\Lambda}U_V(X;A),
\]
whose canonical forest/Mayer activities satisfy
\begin{equation}\label{eq:cell-log}
 \sup_i\sum_{X\ni i}e^{\kappa\ell(X)+|X|}
                         \norm{U_V(X;\cdot)}_\infty\le2e^2\beta.
\end{equation}
The constants are independent of volume. The set $X$ is a set of
marked fluctuation cells; it is not yet a claim of background support.
\end{theorem}

\begin{proof}
BKAR partition-mask covariances are convex combinations of block
principal submatrices, and preserve the stated accretivity and norm
bounds. Allow a further operator perturbation of norm at most $\delta$.
The resulting $M$ has $\Herm M\ge\lambda I/2$, $\norm M\le2\Lambda_0$,
$\Herm M^{-1}\ge V_*^{-1}I$, and
$|\det M|\ge(\lambda/2)^{D|S|}$. Its density modulus is bounded
by $K_G^{|S|}$ times a product of real cell Gaussian densities of
covariance $V_*I$. Hence
\[
 \left|\int p_M(x_S)\prod_{i\in S}q_i(A,x_i)\,dx_S\right|
 \le K_G^{|S|}\prod_{i\in S}
          \E_{V_*I}(\mathbf1_{\{|G_i|>R\}}+v)
 \le\beta^{|S|}.
\]
The tail estimate is exponential Markov with parameter $1/(4V_*)$.
The same domination establishes background holomorphy.

For a tree $T$ with degrees $d_i$, independent covariance-edge
parameter changes of radii
$r_{ij}=\delta/[\max(d_i,d_j)J_{ij}]$ have every block row and
column perturbation bounded by $\delta$. Cauchy's formula and
$d^d\le e^d d!$ give
\[
 |\partial_Tf_S|\le\beta^{|S|}(e^2/\delta)^{|S|-1}
                          \prod_i d_i!\prod_{ij\in T}J_{ij}.
\]
The multipliers are independent of the covariance interpolation
parameters. At a disconnected forest mask their product expectation
factorizes over the components. BKAR thus yields a hard-core gas of
disjoint marked sets with the exact tree activities.

For $n\ge2$ the degree sum is
\[
 \sum_{T\text{ on }[n]}\prod_i d_i!
 =(n-2)!\binom{3n-3}{n-2},\qquad
 \frac{1}{(n-1)!}\sum_T\prod_i d_i!\le8^{n-1}.
\]
Summing leaves with the weighted $J$ norm proves
\[
 \sup_i\sum_{S\ni i}\norm{\zeta(S)}_\infty
             e^{\kappa\ell(S)+\alpha|S|}
 \le\frac{\beta e^\alpha}{1-B_0\beta e^\alpha}.
\]
At $\alpha=2$, the stated ceiling implies the Koteck\'y--Preiss
criterion for activities weighted by $e^{\kappa\ell(S)+|S|}$
and size function $|S|$. The rooted Mayer bound, grouped by the
union of its overlapping marked polymers, proves \eqref{eq:cell-log}.
This is precisely the proof of Notebook 03 with the marked-tail
parameter replaced by $K_G(p_R+v)$.
\end{proof}

\begin{corollary}[The interaction changes the logarithm by $O(v)$]
If $\tau:=K_Gp_R\le\beta_*/4$ and $K_Gv\le\beta_*/4$, the same
canonical activities for $V$ and for $V=0$ obey
\begin{equation}\label{eq:interaction-difference}
 \sup_i\sum_{X\ni i}e^{\kappa\ell(X)+|X|}
       \norm{U_V(X)-U_0(X)}_\infty\le8e^2K_Gv.
\end{equation}
\end{corollary}

\begin{proof}
Interpolate the bounded multipliers by
$q_i(t)=-\mathbf1_{{\rm bad},i}+t\chi_i(e^{V_i}-1)$.
For $v>0$, the disk $|t|\le\rho=\beta_*/(2K_Gv)\ge2$ has
parameter at most $3\beta_*/4$. The activity vector is analytic
there in the weighted norm and bounded by $M=2e^2\beta_*$.
Banach-valued Cauchy bounds for its Taylor coefficients give
\[
 \norm{U(1)-U(0)}\le\frac{M}{\rho-1}
                  \le8e^2K_Gv.
\]
The $v=0$ case is immediate. Both logarithms use the expansion's
branches; no lower bound for a cutoff probability is used.
\end{proof}

\subsection{The strict-background-local version}

Retain \emph{all} hypotheses and smallness conditions of Notebook 04,
including its fixed background-independent baseline, local precision
decomposition, product domain, and constants. Require additionally
that $V_i$ reads only background coordinates in cell $i$.
In the rooted support algebra for a marked set $S$, its extra multiplier
has support $\{i\}\subset S$, zero added support length, and norm at
most $v$. Consequently the multiplier product norm is bounded pointwise
by $\prod_{i\in S}(\mathbf1_{{\rm bad},i}+v)$.

The marginal-density estimate of that notebook is
\[
 \norm{p_M(x_S)}_{\mathcal B_S}
 \le\widehat K_G^{|S|}p_{2V_*I}(x_S),\qquad
 \widehat K_G=(8V_*/\lambda)^{D/2}.
\]
Its proof therefore goes through with the replacement
\[
 \widehat\beta=\widehat K_G
        [2^{D/2}e^{-R^2/(8V_*)}+v].
\]
Use that notebook's $\delta=\lambda/4$, edge envelope, $B_0$, and
ceiling $\widehat\beta\le[2e^2\max(1,B_0)]^{-1}$.
Forest factorization is still valid coefficientwise in the local
source variables, and compatibility is disjointness of the marks,
not of their added background supports. Its connected-hull construction
then gives activities $E_{\Lambda,H}(A|_H)$ with
\[
 \sup_b\sum_{H\ni b}e^{a_*|H|}\norm{E_{\Lambda,H}}_\infty
 \le2e^2S_{\kappa/2}\widehat\beta,
 \qquad a_*=\min(1/2,\kappa/4),
\]
where $S_{\kappa/2}=\sum_{z\in\mathbb Z^d}e^{-(\kappa/2)|z|}$.
Under the analogous quarter budgets, the difference from the
Gaussian-cutoff activities is at most
$8e^2S_{\kappa/2}\widehat K_Gv$ in this norm.

This statement would be false without controlling the interaction's
background support. An interaction in cell $i$ that reads an arbitrarily
remote background variable cannot be assigned support $\{i\}$.
Nor does strict dependence on the varying background imply consistency
between volumes: the finite-volume coefficients may still depend on
the fixed whole-volume baseline. The symmetry and domain caveats of
Notebook 04 remain in force.

\section{The actual enlarged domain and the remaining interaction theorem}

Write $L_g=\log g^{-2}$, $R=A_1L_g^{p_0}$ and
$P=C_1L_g^{p_1}$ with $p_1<p_0$. For every fixed $q$,
\[
 gR^q=A_1^qe^{-L_g/2}L_g^{qp_0}\longrightarrow0.
\]
Thus a fixed physical-field analytic radius permits the larger real
cutoff, after separately enforcing $gR<r_0$. It does not require
$R\le P$, which in fact fails for small $g$. A hypothesis $gP<r_0$
alone is insufficient; a domain only of unscaled radius $P$ cannot
be enlarged by Schwarz's lemma. After the primitive constants are
fixed, choose the coupling ceiling last to meet the cluster budget
and any prescribed absolute error. The constant in $O(gR^3)$ is
independent of that later target; the ceiling need not be.

When $p_0>1/2$, the Gaussian tail $Ke^{-cR^2}$ is smaller than every
fixed $gL_g^q$ scale for sufficiently small $g$, since
$-cA_1^2L_g^{2p_0}+L_g/2-q\log L_g\to-\infty$.

The deposit already contains the appropriate larger-radius calculation.
Raw II, \texttt{5.2:lem-complex-localisation}, proof Step 3,
lines 11178--11195, bounds cubic pieces by
$C(M,\kappa_1)g_k(T^{\C})^3e^{-\kappa_Vd_k(Y)}$, treats merely
vanishing pieces by a $g_kT^{\C}$ factor, and chooses the coupling
ceiling after the activity target. Its domain calculation uses
\[
 \theta_k^{\C}=K_Rg_kT^{\C}/\alpha_{*,k},\qquad
 \mathsf w^{\C}=1/\theta_k^{\C}-1,
 \qquad(1+\mathsf w^{\C})K_Rg_kT^{\C}=\alpha_{*,k}.
\]
It invokes physical response containment; it does not claim that
the larger logarithmic radius fits inside the smaller one. The
remaining verification is its imported response theorem, complex
background chart, admitted histories, and summable localized
coefficient bounds. Raw II's later box-exponent lemma also assumes
a coupling-flow hypothesis, so it cannot independently prove that
same hypothesis.

\subsection{Why overlapping interactions need more than the cell theorem}

For a potential $\sum_YV_Y(A,B_Y)$ with overlapping cell sets, define
$\chi_Y=\prod_{i\in Y}\chi_i$ and
$w_Y=\chi_Y(e^{V_Y}-1)$, setting $w_Y=0$ outside its cutoff.
There is an exact identity
\[
 \left(\prod_i\chi_i\right)e^{\sum_YV_Y}
 =\left(\prod_i\chi_i\right)\prod_Y(1+w_Y),
 \qquad \norm{w_Y}_\infty\le e^{u_Y}-1.
\]
If any cutoff fails, both sides vanish; otherwise the identity is
the product of exponentials. A bad mark in $Y$ also annihilates a
product containing $w_Y$.

These facts do not imply Theorem~\ref{thm:cell-interaction} for
overlapping interactions. A covariance mask fails to factorize across
a partition cut by an interaction hyperedge. Grouping overlapping
interactions into a molecule $S$ makes a Gaussian block of dimension
$D|S|$. A crude edge bound can then pay $|S|$ at every incident edge,
hence $|S|^{d_S}$ at degree $d_S$. The degree factorials from Cauchy's
formula can consume the factorial savings in tree enumeration.
One exponential factor in $|S|$ does not control unrestricted powers
of $|S|$ in that argument. Merely calling these molecules ``cells''
therefore does not prove convergence.

A natural prospective hypothesis is a small weighted norm
\[
 \sup_i\sum_{Y\ni i}e^{a|Y|+\kappa\ell(Y)}(e^{u_Y}-1),
\]
with enough reserve for dimensional factors. No theorem that this
hypothesis alone completes the proposed molecule argument is claimed
here. One needs a valid allocation of forest edges to fixed cells,
or an interaction-specific analytic expansion. CMP116, pages 10--11
and 16--17, uses a field-dependent quadratic-form decomposition and
a separate Gaussian/determinant resummation. The next precise model
target is to verify that theorem on the enlarged physical domain,
with its actual complex operators and history norm.

\section{What these results contribute to the construction}

The flat gauge bridge uses the genuine gauge penalty and both genuine
constraints. It can transfer a matched axial estimate to the Gaussian
needed for the cutoff calculation, with explicit coordinate constants.
Companion Notebook 06 now supplies a direct flat axial gap and the
complementary energy upper bound, using curvature contraction and a
compact right inverse to averaging. This replaces the imported flat
spectral input for the whole-lattice or periodic quadratic model;
the actual smooth minimizer/constraint identification and complex
operator estimates remain explicit.
The physical-field estimates explain how the larger cutoff can be
handled and identify the linear term that a cubic-only argument misses.
The fixed-cell theorem closes a complete, limited interaction module.

For the scalar quadratic cutoff coefficient, Raw I, lines 3780--3837,
extracts the carrier Hessian at $(U,J)=(1,0)$ followed by normalized
responses and a moment sum. A genuinely flat baseline can therefore
be sufficient for the direct-Hessian route of Notebook 03. Choosing
a different baseline $C_0(U_{\rm ref})$ at every real reference only
localizes perturbation dependence; it leaves arbitrary dependence on
$U_{\rm ref}$ and proves no full-background locality.

Still required are uniform covariance/response bounds for the actual
independent-current family, the overlapping interaction estimate, and
their match to the effective-action induction. Beyond this local RG
step, continuum compactness, a sufficiently rich observable hierarchy,
rotation restoration, Yang--Mills identification, nontriviality, and
a mass gap remain separate obligations listed in Notebook 01.

\begin{thebibliography}{9}
\bibitem{balaban99} T. Ba\l aban,
\emph{Propagators for lattice gauge theories in a background field},
Commun. Math. Phys. 99 (1985), 389--434.
\href{https://doi.org/10.1007/BF01240355}{Publisher record}.
\bibitem{balaban109} T. Ba\l aban,
\emph{Renormalization group approach to lattice gauge field theories I},
Commun. Math. Phys. 109 (1987), 249--301.
\href{https://scispace.com/pdf/renormalization-group-approach-to-lattice-gauge-field-3ygai0vxt1.pdf}
{Copy of the original article}.
\bibitem{balaban116} T. Ba\l aban,
\emph{Renormalization group approach to lattice gauge field theories II.
Cluster expansions}, Commun. Math. Phys. 116 (1988), 1--22.
\href{https://scispace.com/pdf/renormalization-group-approach-to-lattice-gauge-field-4penm4k7hc.pdf}
{Copy of the original article}.
\bibitem{bkar} A. Abdesselam,
\href{https://mabdesselam.github.io/BKAR.pdf}{Author-hosted notes on the BKAR formula}.
\bibitem{kp} R. Koteck\'y and D. Preiss,
\emph{Cluster expansion for abstract polymer models},
Commun. Math. Phys. 103 (1986), 491--498.
\href{https://doi.org/10.1007/BF01211762}{Publisher record}.
\end{thebibliography}
\end{document}
